\section{Conclusion}
This paper presents the measurement of charged hadron multiplicity per unit pseudorapidity, \dndeta, using field-off data from Run 2024 in \auau collisions at $\sqrtsnn=200\gev$, collected with the sPHENIX detector. Results are reported as a function of pseudorapidity, $\eta$, within $|\eta|<1.1$ across multiple centrality intervals. In the $3\%$ most central collisions, \dndeta for $|\eta|<0.3$ is measured to be $723.4\pm45.3$. The \dndeta strongly increases with collision centrality, and the normalised \dndeta per participant pair exhibits a mild rise with the number of participating nucleons, \npart. 
% This observable is further compared to predictions from the \hijing, \epos, and \ampt Monte Carlo heavy-ion event generators, with results indicating that the \sphenix measurement is most compatible with \hijing in both magnitude and \npart dependence. 
The \sphenix measurement, which combines two analysis methods, provides full $2\pi$ azimuthal coverage at mid-rapidity and presents a characterization of charged particle multiplicity as a function of $\eta$. These results are in agreement with previous measurements from \phobos, \phenix, and \brahms, further supporting the broader sPHENIX physics program.

% The analysis also serves an essential commissioning purpose, by demonstrating the capabilities of several new detector components and their agreement with established physics results, which will further enable the broader \sphenix physics program.